% Options for packages loaded elsewhere
\PassOptionsToPackage{unicode}{hyperref}
\PassOptionsToPackage{hyphens}{url}
\PassOptionsToPackage{dvipsnames,svgnames,x11names}{xcolor}
\documentclass[
  10pt,
]{article}
\usepackage{xcolor}
\usepackage[margin=1cm,top=1cm,bottom=1cm,left=1cm,right=1cm,includeheadfoot]{geometry}
\usepackage{amsmath,amssymb}
\setcounter{secnumdepth}{3}
\usepackage{iftex}
\ifPDFTeX
  \usepackage[T1]{fontenc}
  \usepackage[utf8]{inputenc}
  \usepackage{textcomp} % provide euro and other symbols
\else % if luatex or xetex
  \usepackage{unicode-math} % this also loads fontspec
  \defaultfontfeatures{Scale=MatchLowercase}
  \defaultfontfeatures[\rmfamily]{Ligatures=TeX,Scale=1}
\fi
\usepackage{lmodern}
\ifPDFTeX\else
  % xetex/luatex font selection
  \setmainfont[]{DejaVu Serif}
  \setmonofont[]{DejaVu Sans Mono}
\fi
% Use upquote if available, for straight quotes in verbatim environments
\IfFileExists{upquote.sty}{\usepackage{upquote}}{}
\IfFileExists{microtype.sty}{% use microtype if available
  \usepackage[]{microtype}
  \UseMicrotypeSet[protrusion]{basicmath} % disable protrusion for tt fonts
}{}
\usepackage{setspace}
\makeatletter
\@ifundefined{KOMAClassName}{% if non-KOMA class
  \IfFileExists{parskip.sty}{%
    \usepackage{parskip}
  }{% else
    \setlength{\parindent}{0pt}
    \setlength{\parskip}{6pt plus 2pt minus 1pt}}
}{% if KOMA class
  \KOMAoptions{parskip=half}}
\makeatother
\usepackage{listings}
\newcommand{\passthrough}[1]{#1}
\lstset{defaultdialect=[5.3]Lua}
\lstset{defaultdialect=[x86masm]Assembler}
\usepackage{graphicx}
\makeatletter
\newsavebox\pandoc@box
\newcommand*\pandocbounded[1]{% scales image to fit in text height/width
  \sbox\pandoc@box{#1}%
  \Gscale@div\@tempa{\textheight}{\dimexpr\ht\pandoc@box+\dp\pandoc@box\relax}%
  \Gscale@div\@tempb{\linewidth}{\wd\pandoc@box}%
  \ifdim\@tempb\p@<\@tempa\p@\let\@tempa\@tempb\fi% select the smaller of both
  \ifdim\@tempa\p@<\p@\scalebox{\@tempa}{\usebox\pandoc@box}%
  \else\usebox{\pandoc@box}%
  \fi%
}
% Set default figure placement to htbp
\def\fps@figure{htbp}
\makeatother
\setlength{\emergencystretch}{3em} % prevent overfull lines
\providecommand{\tightlist}{%
  \setlength{\itemsep}{0pt}\setlength{\parskip}{0pt}}
% Enable graphics inclusion and ensure figure numbering works
\usepackage{graphicx}
\renewcommand{\figurename}{Figure}

% Configure fonts for Unicode support with fallbacks
\usepackage{newunicodechar}
\newunicodechar{⁴}{\textsuperscript{4}}
\newunicodechar{₄}{\textsubscript{4}}

% Math notation macros
\newcommand{\norm}[1]{\lVert #1\rVert}
\newcommand{\abs}[1]{\lvert #1\rvert}

% Enhanced code block styling for better contrast and readability
\usepackage{fancyvrb}
\usepackage{xcolor}
\usepackage{listings}

% Define custom colors for code blocks
\definecolor{codebg}{RGB}{245, 245, 245}      % Light gray background
\definecolor{codeborder}{RGB}{200, 200, 200}  % Medium gray border
\definecolor{codefg}{RGB}{50, 50, 50}         % Dark gray text

% Configure Verbatim environment for inline code
\DefineVerbatimEnvironment{Verbatim}{Verbatim}{%
    fontsize=\small,
    frame=single,
    framerule=0.5pt,
    framesep=3pt,
    rulecolor=\color{codeborder},
    bgcolor=\color{codebg},
    fgcolor=\color{codefg}
}

% Configure code block styling
\DefineVerbatimEnvironment{Highlighting}{Verbatim}{%
    fontsize=\footnotesize,
    frame=single,
    framerule=0.5pt,
    framesep=5pt,
    rulecolor=\color{codeborder},
    bgcolor=\color{codebg},
    fgcolor=\color{codefg}
}

% Style inline code with \texttt
\renewcommand{\texttt}[1]{%
    \colorbox{codebg}{\color{codefg}\ttfamily #1}%
}

% Configure listings package for code blocks
\lstset{
    backgroundcolor=\color{codebg},
    basicstyle=\footnotesize\ttfamily\color{codefg},
    breakatwhitespace=false,
    breaklines=true,
    captionpos=b,
    commentstyle=\color{codefg},
    deletekeywords={...},
    escapeinside={\%*}{*)},
    extendedchars=true,
    frame=single,
    framerule=0.5pt,
    framesep=5pt,
    keepspaces=true,
    keywordstyle=\color{codefg},
    language=Python,
    morekeywords={*,...},
    numbers=left,
    numbersep=5pt,
    numberstyle=\tiny\color{codefg},
    rulecolor=\color{codeborder},
    showspaces=false,
    showstringspaces=false,
    showtabs=false,
    stepnumber=1,
    stringstyle=\color{codefg},
    tabsize=2,
    title=\lstname
}

% Override any Pandoc default lstset configurations
\AtBeginDocument{
    \lstset{
        backgroundcolor=\color{codebg},
        basicstyle=\footnotesize\ttfamily\color{codefg},
        frame=single,
        framerule=0.5pt,
        framesep=5pt,
        rulecolor=\color{codeborder},
        numbers=left,
        numbersep=5pt,
        numberstyle=\tiny\color{codefg}
    }
}

% Configure hyperref colors consistently
\AtBeginDocument{
% Override pandoc's hidelinks setting with consistent options
\hypersetup{
    colorlinks=true,
    allcolors=red,
    linkcolor=red,
    urlcolor=red,
    citecolor=red,
    filecolor=red,
    menucolor=red,
    linktoc=all
}
}

% Simple page break support for document structure
\usepackage{bookmark}
\IfFileExists{xurl.sty}{\usepackage{xurl}}{} % add URL line breaks if available
\urlstyle{same}
% fallback for those not using the hyperref driver hyperxmp:
\makeatletter
\@ifundefined{xmpquote}{\newcommand{\xmpquote}[1]{#1}}{}
\makeatother
\hypersetup{
  colorlinks=true,
  linkcolor={red},
  filecolor={red},
  citecolor={red},
  urlcolor={red},
  pdfcreator={LaTeX via pandoc}}

\title{Lattice Visualization Gallery}
\author{Daniel Ari Friedman\\ ORCID: 0000-0001-6232-9096\\ Email: daniel@activeinference.institute\\ DOI: 10.5281/zenodo.16887791}
\date{October 07, 2026}

\begin{document}
\maketitle

{
\hypersetup{linkcolor=black}
\setcounter{tocdepth}{3}
\tableofcontents
}
\setstretch{1.0}
\section{Lattice Visualization
Gallery}\label{lattice-visualization-gallery}

\subsection{Overview}\label{overview}

This section is the figure surface of the lattice layer: every rendering
primitive used by the IVM chapters lives in
\passthrough{\lstinline!vis\_lattice.py!}
(\passthrough{\lstinline!shell\_scatter!},
\passthrough{\lstinline!field\_slice!},
\passthrough{\lstinline!dynamics\_strip!},
\passthrough{\lstinline!gallery!}). Each primitive receives its
matplotlib axes explicitly, so panels compose inside caller-owned
figures; only \passthrough{\lstinline!gallery!} creates figures (three,
written as PNGs). Nothing renders at import time, and all randomness is
drawn from a seeded \passthrough{\lstinline!numpy.random.default\_rng!},
so the gallery is deterministic. The command-line entry is
\passthrough{\lstinline!quadmath/scripts/lattice\_gallery.py!} --- a
thin orchestrator that sets a headless backend and a fixed seed,
delegates to \passthrough{\lstinline!gallery!} (thin-orchestrator
contract of \passthrough{\lstinline!quadmath/scripts/AGENTS.md!}), and
prints each written path on its own line: those stdout lines are the
\passthrough{\lstinline!make\_all\_figures!} manifest contract.

\subsection{Frequency shells}\label{frequency-shells}

\passthrough{\lstinline!shell\_scatter(ax, sites, k)!} embeds the sites
of shell \passthrough{\lstinline!k!} (from
:func:\passthrough{\lstinline!omni\_numbering.generate\_shell!}, which
enumerates the \(10k^2 + 2\) sites of shell \(k\); see
\href{13_lattice_tooling.md}{Lattice Tooling}) with
\passthrough{\lstinline!to\_xyz!} under
\passthrough{\lstinline!DEFAULT\_EMBEDDING!} and scatters them in 3D.
With \passthrough{\lstinline!axis\_hints!} it overlays four dashed rays
from the origin toward the embedding matrix columns --- the images of
the quadray basis directions \(A, B, C, D\), which are the vertices of a
regular tetrahedron (pairwise dot product \(-1\), edge length
\(2\sqrt{2}\); see \href{03_quadray_methods.md}{Quadray Methods}). The
hints make the tetrahedral axis frame of the lattice readable in print.

\begin{figure}
\centering
\pandocbounded{\includegraphics[keepaspectratio,alt={Frequency shells of the omnidirectional close packing. Shell 1 (12 sites) and shell 2 (42 sites) of the IVM lattice embedded via DEFAULT\_EMBEDDING and scattered by shell\_scatter as 3D point clouds in the x, y, z embedding coordinates; dashed gray rays mark the four tetrahedral quadray axes, labeled A, B, C, D, pointing toward the embedding-matrix columns. Reproduced with uv run python quadmath/scripts/lattice\_gallery.py (fixed seed 12).}]{../output/figures/vis_gallery_shell.png}}
\caption{\textbf{Frequency shells of the omnidirectional close packing.}
Shell 1 (12 sites) and shell 2 (42 sites) of the IVM lattice embedded
via \protect\passthrough{\lstinline!DEFAULT\_EMBEDDING!} and scattered
by \protect\passthrough{\lstinline!shell\_scatter!} as 3D point clouds
in the \(x, y, z\) embedding coordinates; dashed gray rays mark the four
tetrahedral quadray axes, labeled A, B, C, D, pointing toward the
embedding-matrix columns. Reproduced with
\protect\passthrough{\lstinline!uv run python quadmath/scripts/lattice\_gallery.py!}
(fixed seed 12).}
\end{figure}

\subsection{Lattice-plane slices}\label{lattice-plane-slices}

\passthrough{\lstinline!field\_slice(ax, field, sites, plane)!} renders
a scalar \passthrough{\lstinline!ivm\_field.IVMField!} as a heatmap
restricted to a lattice plane. A plane is parametrized by an origin
\(q_0\) and two independent integer translation vectors \(u\), \(v\) ---
by default two of the twelve IVM neighbor moves, both permutations of
\((2,1,1,0)\) --- and its points are

\begin{equation}
\label{eq:vis-plane}
q(i, j) \;=\; \operatorname{normalize}\!\big(q_0 + i\,u + j\,v\big),
\qquad (i, j) \in \mathbb{Z}^2 ,
\end{equation}

with \passthrough{\lstinline!normalize!} the quadray canonical
representative. Because each neighbor move has component sum \(4\), the
sum stays divisible by \(4\) along the whole plane, so every \(q(i,j)\)
is again a lattice site. Two coordinate facts make the rendering exact:
the embedding is linear, so the xyz image is the planar set
\(t(q_0) + i\,t(u) + j\,t(v)\); and the shift \((1,1,1,1)\) lies in the
embedding kernel, so plane membership is tested in integer quadray space
--- a site lies on the plane of \eqref{eq:vis-plane} iff
\(q - q_0 = i\,u + j\,v + m\,(1,1,1,1)\) for integers \(i, j, m\).
\passthrough{\lstinline!field\_slice!} decides membership by an exact
least-squares solve followed by an integer verification, so no site is
ever misassigned by floating point. Sites off the plane are ignored,
plane cells outside the field's ball are masked, and the heatmap axes
are the integer plane indices \(i\) (steps along \(u\)) and \(j\) (steps
along \(v\)).

The gallery figure learns a field on the radius-3 ball (147 sites) from
noisy observations of a linear synthetic truth (25\% of sites observed,
noise level \(0.1\), regularization \passthrough{\lstinline!lam=0.01!},
kernel width \passthrough{\lstinline!0.5!}) and slices it through the
default plane \(\big((2,1,1,0), (1,2,1,0)\big)\); the learned surface
follows the planar lattice, with masked tiles where the ball ends. Field
learning itself is the subject of \href{11_ivm_field_learning.md}{Static
IVM Field Learning}.

\begin{figure}
\centering
\pandocbounded{\includegraphics[keepaspectratio,alt={Learned scalar field sliced along a lattice plane. Heatmap of a Laplacian-regularized IVMField over the radius-3 ball, restricted by field\_slice to the plane spanned by the neighbor moves u = (2,1,1,0) and v = (1,2,1,0): a viridis heatmap whose axes are the integer plane indices i (steps along u) and j (steps along v), with the colorbar reporting the field value and light-gray masked tiles marking plane cells outside the ball. Reproduced with uv run python quadmath/scripts/lattice\_gallery.py (fixed seed 12).}]{../output/figures/vis_gallery_field.png}}
\caption{\textbf{Learned scalar field sliced along a lattice plane.}
Heatmap of a Laplacian-regularized
\protect\passthrough{\lstinline!IVMField!} over the radius-3 ball,
restricted by \protect\passthrough{\lstinline!field\_slice!} to the
plane spanned by the neighbor moves u = (2,1,1,0) and v = (1,2,1,0): a
\protect\passthrough{\lstinline!viridis!} heatmap whose axes are the
integer plane indices \(i\) (steps along u) and \(j\) (steps along v),
with the colorbar reporting the field value and light-gray masked tiles
marking plane cells outside the ball. Reproduced with
\protect\passthrough{\lstinline!uv run python quadmath/scripts/lattice\_gallery.py!}
(fixed seed 12).}
\end{figure}

\subsection{Dynamics strips}\label{dynamics-strips}

\passthrough{\lstinline!dynamics\_strip(axs, trajectory, t\_indices)!}
renders selected snapshots of an
\passthrough{\lstinline!ivm\_dynamics.simulate!} trajectory (see
\href{12_ivm_dynamics.md}{Dynamics and Learning on the IVM Lattice}) as
a strip of 3D scatter panels over the embedded lattice. All panels share
one symmetric color scale \([-v_{\max}, v_{\max}]\) with
\(v_{\max} = \max_t |u_t|\) over the selected snapshots, so evolution is
directly comparable across panels. The gallery strip shows heat
diffusion (\(\alpha = 0.45\), horizon \(T = 20\)) on the radius-3
lattice (13 sites) at \(t = 0, 10, 20\): the seeded random initial field
relaxes toward its lattice average --- the visually flat mid panel and
right panel are the \(L_2\)-monotone decay of the heat lemma in action.

\begin{figure}
\centering
\pandocbounded{\includegraphics[keepaspectratio,alt={Heat-diffusion evolution strip. Snapshots of a seeded heat run at t = 0, 10, 20 on the radius-3 IVM lattice (13 sites), rendered by dynamics\_strip as 3D scatter panels on one shared symmetric coolwarm scale {[}-v\_\{\textbackslash max\}, v\_\{\textbackslash max\}{]} with v\_\{\textbackslash max\} = \textbackslash max\_t \textbar u\_t\textbar{} over the selected snapshots; the field visibly relaxes as the sum of squares decays monotonically. Reproduced with uv run python quadmath/scripts/lattice\_gallery.py (fixed seed 12).}]{../output/figures/vis_gallery_dynamics.png}}
\caption{\textbf{Heat-diffusion evolution strip.} Snapshots of a seeded
heat run at t = 0, 10, 20 on the radius-3 IVM lattice (13 sites),
rendered by \protect\passthrough{\lstinline!dynamics\_strip!} as 3D
scatter panels on one shared symmetric
\protect\passthrough{\lstinline!coolwarm!} scale
\([-v_{\max}, v_{\max}]\) with \(v_{\max} = \max_t |u_t|\) over the
selected snapshots; the field visibly relaxes as the sum of squares
decays monotonically. Reproduced with
\protect\passthrough{\lstinline!uv run python quadmath/scripts/lattice\_gallery.py!}
(fixed seed 12).}
\end{figure}

\subsection{Deterministic animations and gallery
plots}\label{sec:animation_frames}

Where the primitives above compose panels inside caller-owned figures,
the module \passthrough{\lstinline!src/quadmath/viz/animations.py!}
renders frame-based animations that own their pixels: each renderer
returns a list of \passthrough{\lstinline!Frame!} objects --- the frozen
\passthrough{\lstinline!dataclass!} \passthrough{\lstinline!Frame!} in
that file pairs a 2D \passthrough{\lstinline!numpy!} array with a
human-readable \passthrough{\lstinline!title!}, and its
\passthrough{\lstinline!\_\_post\_init\_\_!} validation admits only
\passthrough{\lstinline!uint8!} arrays (raw gray levels) or float arrays
with every value in \([0, 1]\), rejecting anything that is not 2D, is of
another dtype, or drifts outside the unit interval. No matplotlib is
involved at frame level and no RNG or wall-clock input exists (except
the explicit \passthrough{\lstinline!seed!} of
\passthrough{\lstinline!diffusion\_frames!}), so identical arguments
yield byte-identical frames. All three renderers share one fixed-camera
convention: sites are embedded in \(R^3\) via
\passthrough{\lstinline!DEFAULT\_EMBEDDING!} and viewed by an
orthographic camera looking down \(+z\) --- each point \((x,y,z)\)
projects to \((x,y)\) on a \passthrough{\lstinline!GRID\_SIZE!} \(=48\)
square grid whose extent is \([-8, 8]\) (module constant
\passthrough{\lstinline!\_HALF\_EXTENT!}, sized to cover a radius-3 IVM
ball of embedded norm \(6\) with margin even while pulsing), the
vertical axis flipped so \(+y\) points up, colliding points resolved by
maximum brightness, and out-of-window points clamped to the border
pixel.

\passthrough{\lstinline!simplex\_frames(q0, q1, n=16)!} animates a
quaternion-slerp rotation of the radius-1 ball: the two unit quaternions
\(q_0\), \(q_1\) (validated to four components of unit norm) are
interpolated on the shortest arc --- sign-flipping \(q_1\) when their
dot product is negative, using the Shoemaker spherical formula
\(q(t) = \sin((1-t)\theta)\,q_0 / \sin\theta + \sin(t\theta)\,q_1 / \sin\theta\),
and falling back to normalized linear interpolation when the inputs are
nearly parallel --- and at each \(t = i/(n-1)\) the 13 lattice sites of
\passthrough{\lstinline!ball\_sites!} radius 1 are rotated by \(q(t)\),
projected onto the 48×48 grid, and shaded by radial shell, with origin
brightness \(1.0\) and the outermost shell never dimmer than \(0.25\) so
nothing renders invisible.

\passthrough{\lstinline!lattice\_frames(shells=3, n=12)!} renders a
pulsing IVM lattice ball: the radius-\passthrough{\lstinline!shells!}
ball from \passthrough{\lstinline!ball\_sites!} (147 sites at the
default radius 3) has its embedded coordinates scaled per frame by the
deterministic pulse \(1 + 0.25\sin(2\pi i/n)\) over one full
growth-and-contraction cycle, while brightness per site stays the fixed
radial-shell shading --- brightest at the center, decaying linearly in
shell index, outermost shell at \(0.25\) --- so the pulse animates
geometry, not shading.

\passthrough{\lstinline!diffusion\_frames(n\_steps=12, seed=0)!} seeds
one-hot heat on the radius-3 ball and diffuses it over the IVM neighbor
graph: the 147 ball sites become a graph in which two sites are adjacent
iff their difference normalizes to one of the twelve
\passthrough{\lstinline!IVM\_NEIGHBOR\_STEPS!}, a
\passthrough{\lstinline!numpy.random.default\_rng(seed)!} draw selects
the source site deterministically, and each explicit update applies
\(u \leftarrow u + \alpha\,(\text{mean of graph neighbors} - u)\) with
\(\alpha = 0.25\) (stable below \(0.5\)) and values clipped at zero.
Each frame is normalized by dividing by its (always positive) maximum
heat, giving float values in \([0, 1]\); because heat only spreads to
graph neighbors, the set of heated sites grows monotonically across
frames.

The companion module \passthrough{\lstinline!src/quadmath/viz/plots.py!}
supplies standalone figure builders for diagnostics that do not belong
to any gallery: \passthrough{\lstinline!plot\_loss\_history!} draws a
training loss sequence as a marker line over its iteration index,
\passthrough{\lstinline!plot\_shell\_growth!} plots the IVM shell
cardinalities --- the cuboctahedral numbers \(10k^2+2\) produced by
\passthrough{\lstinline!shell\_cardinalities!} in
\passthrough{\lstinline!src/quadmath/lattice/ivm\_field.py!} --- versus
shell index, \passthrough{\lstinline!plot\_error\_histogram!} histograms
error values with a dashed vertical mean line, and
\passthrough{\lstinline!plot\_lattice\_shell\_3d!} scatters the sites of
one shell (\passthrough{\lstinline!shell\_sites!}, embedded via
\passthrough{\lstinline!DEFAULT\_EMBEDDING!}) in 3D with equal-aspect
axes. Unlike the axes-composing primitives of
\passthrough{\lstinline!vis\_lattice!}, each owns a fixed-size figure
(6.4×4.8 inches at 160 dpi), applies a fixed style, and on
\passthrough{\lstinline!save=True!} writes its PNG into
\passthrough{\lstinline!quadmath/output/figures/!} via
\passthrough{\lstinline!quadmath.paths.get\_figure\_dir!}, returning the
output path (empty string when the caller keeps the figure). All four
are input-agnostic and seed-free: inputs are plain sequences or shell
indices, styling is fixed, and no randomness or wall-clock time enters,
so runs stay deterministic headless
(\passthrough{\lstinline!MPLBACKEND=Agg!}).

\passthrough{\lstinline!frames\_to\_gif(frames, out\_path, fps=8, scale=8)!}
assembles a frame sequence into an animated GIF: float frames are
quantized to \passthrough{\lstinline!uint8!} gray levels by rounding
\(255\,v\), each frame is upscaled by the integer factor
\passthrough{\lstinline!scale!} with nearest-neighbor resampling, and
the sequence is saved with per-frame duration \(1000\,/\,\texttt{fps}\)
milliseconds and infinite looping. PIL is imported lazily inside the
function (Pillow stays optional at import time), and because the PIL GIF
encoder is deterministic and embeds no timestamps, identical
\passthrough{\lstinline!(frames, fps, scale)!} inputs render
byte-identical GIF files --- the same byte-stability contract the PNG
gallery asserts.

\passthrough{\lstinline!frames\_strip(frames, out\_path=None, labels=None, save=True)!}
renders the print still-counterpart of
\passthrough{\lstinline!frames\_to\_gif!}: a single-row matplotlib strip
with one grayscale \passthrough{\lstinline!imshow!} panel per
\passthrough{\lstinline!Frame!}, figure width scaling with the panel
count while the height stays one panel plus the title band, each panel
pinned to the unit brightness interval
(\passthrough{\lstinline!vmin = 0!}, \passthrough{\lstinline!vmax = 1!};
\passthrough{\lstinline!uint8!} frames are first mapped by \(v/255\), so
both \passthrough{\lstinline!Frame!} dtypes share the GIF's gray-level
convention), per-panel titles taken from
\passthrough{\lstinline!labels!} or, by default, from each frame's own
\passthrough{\lstinline!title!}, and axes switched off with zero
inter-panel spacing so consecutive panels butt together. Matplotlib is
imported lazily inside the function (the module stays matplotlib-free at
import time), and saving follows the \passthrough{\lstinline!plots.py!}
convention: a bare \passthrough{\lstinline!out\_path!} name is written
into \passthrough{\lstinline!quadmath/output/figures/!} via
\passthrough{\lstinline!quadmath.paths.get\_figure\_dir!}, a path
carrying a directory component is used verbatim, and the returned string
is the written path (\passthrough{\lstinline!""!} when
\passthrough{\lstinline!save!} is off or no name is given). An empty
\passthrough{\lstinline!frames!} sequence or a
\passthrough{\lstinline!labels!} length mismatch raises
\passthrough{\lstinline!ValueError!}; the render has no RNG and no
wall-clock input, so identical inputs give byte-identical PNGs --- the
test suite asserts the byte stability, the full \([0, 1]\) grayscale
excursion, the uint8/float equivalence, and the panel-count scaling.

\begin{figure}
\centering
\pandocbounded{\includegraphics[keepaspectratio,alt={Three animation stills in one strip. Single-row frames\_strip composition of one evenly spaced still from each 6-frame animation: left, the pulsing radius-2 IVM ball sampled at the quarter-pulse fraction (lattice\_frames, radial-shell brightness, brightest at the center); center, the quaternion-slerp rotation of the radius-1 ball in the midpoint region of the arc (t = 0.4 on the way from the identity to a 90-degree rotation about z, simplex\_frames); right, the final state of seeded heat diffusion on the radius-3 ball (diffusion\_frames, brightness normalized by the frame's maximum heat). All panels are 48×48 orthographic projections under the fixed +z camera with brightness in {[}0, 1{]}; reproduced with uv run python quadmath/scripts/animation\_stills.py.}]{../output/figures/animation_frames_strip.png}}
\caption{\textbf{Three animation stills in one strip.} Single-row
\protect\passthrough{\lstinline!frames\_strip!} composition of one
evenly spaced still from each 6-frame animation: left, the pulsing
radius-2 IVM ball sampled at the quarter-pulse fraction
(\protect\passthrough{\lstinline!lattice\_frames!}, radial-shell
brightness, brightest at the center); center, the quaternion-slerp
rotation of the radius-1 ball in the midpoint region of the arc
(\(t = 0.4\) on the way from the identity to a 90-degree rotation about
\(z\), \protect\passthrough{\lstinline!simplex\_frames!}); right, the
final state of seeded heat diffusion on the radius-3 ball
(\protect\passthrough{\lstinline!diffusion\_frames!}, brightness
normalized by the frame's maximum heat). All panels are 48×48
orthographic projections under the fixed \(+z\) camera with brightness
in \([0, 1]\); reproduced with
\protect\passthrough{\lstinline!uv run python quadmath/scripts/animation\_stills.py!}.}
\end{figure}

The command-line surface is
\passthrough{\lstinline!quadmath/scripts/animation\_gallery.py!} --- a
thin orchestrator in the
\passthrough{\lstinline!quadmath/scripts/AGENTS.md!} contract that sets
the headless \passthrough{\lstinline!Agg!} backend and renders three
GIFs via \passthrough{\lstinline!frames\_to\_gif!}:
\passthrough{\lstinline!animation\_lattice.gif!} (pulsing radius-3 ball,
12 frames), \passthrough{\lstinline!animation\_simplex.gif!} (slerp from
the identity to a 90-degree rotation about the \(x\) axis, 16 frames),
and \passthrough{\lstinline!animation\_diffusion.gif!} (12 diffusion
steps from seed 0), writing each into
\passthrough{\lstinline!quadmath/output/figures/!} and printing each
path on its own line. The script is registered in
\passthrough{\lstinline!quadmath/scripts/make\_all\_figures.py!}
alongside the other galleries, so the manifest contract of the Overview
applies unchanged: only path lines on stdout, byte-identical re-renders
for identical inputs. The companion
\passthrough{\lstinline!quadmath/scripts/animation\_stills.py!} follows
the same thin-orchestrator contract for print stills: it renders the
three 6-frame sequences of the module ---
\passthrough{\lstinline!lattice\_frames!} at radius 2,
\passthrough{\lstinline!simplex\_frames!} from the identity to a
90-degree rotation about \(z\) with the endpoint quaternion derived from
\passthrough{\lstinline!rotate\_about\_axis!}, and
\passthrough{\lstinline!diffusion\_frames!} at seed 0 --- samples each
at an evenly spaced fraction (quarter-pulse, slerp midpoint region,
final diffusion step), composes the three stills with
\passthrough{\lstinline!frames\_strip!} into
\passthrough{\lstinline!animation\_frames\_strip.png!}, and prints the
written path on its own line.

\subsection{Reproducibility and test
contract}\label{reproducibility-and-test-contract}

\begin{itemize}
\tightlist
\item
  \passthrough{\lstinline!gallery(paths\_out\_dir, seed=12)!} writes
  exactly the three files
  \passthrough{\lstinline!vis\_gallery\_shell.png!},
  \passthrough{\lstinline!vis\_gallery\_field.png!},
  \passthrough{\lstinline!vis\_gallery\_dynamics.png!} (module constant
  \passthrough{\lstinline!GALLERY\_FILES!}, in that order) into
  \passthrough{\lstinline!paths\_out\_dir!}, creating it when missing,
  and returns their paths. Every panel is fully seeded, and the PNGs
  carry no timestamp metadata: the test suite asserts byte-identical
  re-renders for a fixed seed.
\item
  Tests: \passthrough{\lstinline!tests/test\_vis\_lattice.py!} --- 19
  tests, no mocks, deterministic assertions without pixel diffs (artist
  placement on caller-provided axes, exact field values at known lattice
  cells, byte-level reproducibility). Scope with
  \passthrough{\lstinline!PYTEST\_DISABLE\_PLUGIN\_AUTOLOAD=1 uv run coverage run -m pytest tests/test\_vis\_lattice.py -q!}
  and \passthrough{\lstinline!uv run coverage report!} --- 100\%
  statement and branch coverage of
  \passthrough{\lstinline!src/quadmath/viz/vis\_lattice.py!}.
\item
  Tests: \passthrough{\lstinline!tests/unit/viz/test\_animations.py!}
  covers the frame module (\passthrough{\lstinline!Frame!} validation,
  the three renderers, \passthrough{\lstinline!frames\_to\_gif!}) and
  the strip renderer --- \passthrough{\lstinline!frames\_strip!} label
  and emptiness validation, byte-identical re-renders in temporary
  directories, uint8/float equivalence, the full grayscale excursion,
  and panel-count scaling --- headless via
  \passthrough{\lstinline!MPLBACKEND=Agg!} (set in
  \passthrough{\lstinline!tests/conftest.py!}).
\item
  Figure regeneration:
  \passthrough{\lstinline!uv run python quadmath/scripts/lattice\_gallery.py!}
  (the script sets \passthrough{\lstinline!MPLBACKEND=Agg!} itself);
  stdout is exactly the three written paths.
\item
  Still-strip regeneration:
  \passthrough{\lstinline!uv run python quadmath/scripts/animation\_stills.py!}
  (the script sets \passthrough{\lstinline!MPLBACKEND=Agg!} itself);
  stdout is exactly the one written strip path,
  \passthrough{\lstinline!quadmath/output/figures/animation\_frames\_strip.png!}.
\item
  Markdown:
  \passthrough{\lstinline!uv run python quadmath/scripts/validate\_markdown.py!}.
\end{itemize}

\subsection{Cross-references}\label{cross-references}

\begin{itemize}
\tightlist
\item
  Site geometry, normalization, and the embedding:
  \href{03_quadray_methods.md}{Quadray Methods}.
\item
  Shell enumeration and nearest-site search:
  \href{13_lattice_tooling.md}{Lattice Tooling}.
\item
  The sliced learner: \href{11_ivm_field_learning.md}{Static IVM Field
  Learning}.
\item
  The stripped dynamics: \href{12_ivm_dynamics.md}{Dynamics and Learning
  on the IVM Lattice}.
\end{itemize}

\end{document}
